% !PRINTED 153
\weberpage{153}
%<<CONT
$j(\omega)$ and call it simply the \gesperrt{invariant}, and define it thus:
\weq{15}{j(\omega) = 2^8\frac{(1-\varkappa^2\varkappa'^2)^3}{\varkappa^4\varkappa'^4},}
from which we obtain:
\weq{16}{\begin{aligned}
\frac{4.27g_2^3}{G} &= j(\omega), \\
\frac{4.27.27g_3^2}{G} &= j(\omega)-27.64 = \frac{64.(2+\varkappa^2\varkappa'^2)^2(\varkappa'^2-\varkappa^2)^2}{\varkappa^4\varkappa'^4}.
\end{aligned}}
Regarded as a function of $\varkappa^2$, the function $j(\omega)$ has the property of
remaining unchanged if in place of $\varkappa^2$ one of the six values
\[ \varkappa^2, \quad \varkappa'^2, \quad \frac{1}{\varkappa^2}, \quad \frac{1}{\varkappa'^2}, \quad -\frac{\varkappa^2}{\varkappa'^2}, \quad -\frac{\varkappa'^2}{\varkappa^2} \]
is substituted.

If we form the differential quotient of the function $\wp(u)$ according to (7), we obtain
\weq{17}{\wp'(u) = -\frac{16K^3}{\omega_1^3}\frac{\mathrm{cn}\,v\,\mathrm{dn}\,v}{\mathrm{sn}^3v} = -2\frac{\sigma_{00}(u)\sigma_{10}(u)\sigma_{01}(u)}{\sigma(u)^3}}
and from this, by (3):
\weq{18}{\begin{aligned}
\wp'(u)^2 &= 4[\wp(u)-e_1][\wp(u)-e_2][\wp(u)-e_3] \\
&= 4\wp(u)^3-g_2\wp(u)-g_3,
\end{aligned}}
or finally
\weq{19}{du = \frac{d\wp}{\sqrt{4\wp^3-g_2\wp-g_3}},}
from which one sees that the function $\wp(u)$ stands in the same relation to the
\gesperrt{Weierstrass} normal form of the elliptic differential as the function
$\mathrm{sn}\,v$ does to the \gesperrt{Legendre} one.

\begin{center}
\textbf{§\,47. \gesperrt{The elliptic transcendents of the second kind.}}
\end{center}

\gesperrt{Jacobi} introduced as the transcendent of the second kind the function
\weq{1}{Z(v) = \frac{\Theta'(v)}{\Theta(v)} = \frac{d\log\Theta(v)}{dv}}
The relation of this function to the elliptic integrals of the second kind follows from
formula (16) of §\,23; and it may be remarked at once that quite similar considerations
could be attached to the formulae (15), (17), (18) there, which, however, do not lead to
essentially new results.
